\section{\faWrench\ Multimodalità\\\small{Speech-to-text}} % (fold)
\label{sec:spring-stt}
%
\begin{frame}[t,fragile] \frametitle{Multimodalità}
    \framesubtitle{Speech-to-text}
    {\footnotesize
    \begin{itemize}[leftmargin=10pt,align=right]
        \onslide<1->\item[\alert{\faArrowCircleRight}] Processo di produzione testuale da audio del parlato
        \onslide<2->\item[\alert{\faArrowCircleRight}] Compatibilità Spring AI (al 11/2025)
        \begin{itemize}[leftmargin=10pt,align=right]
            \onslide<3->\item[\alert{\faArrowCircleRight}] OpenAI API
            \item[\alert{\faArrowCircleRight}] Azure OpenAI
            \item[\alert{\faArrowCircleRight}] OpenAI-\textit{compatible} API (limitata)
        \end{itemize}
    \end{itemize}
    }
\end{frame}    
%
\begin{frame}[t,fragile] \frametitle{Multimodalità}
    \framesubtitle{Speech-to-text: caso OpenAI}
    {\footnotesize
        \begin{block}{Dipendenze applicativo}
			{\tiny\inputminted{xml}{code/pom.xml}}
    	\end{block}
        \begin{block}{Configurazione applicativo}
			{\tiny\inputminted{yaml}{code/application.yml}}
    	\end{block}       
    }
\end{frame}
%
\begin{frame}[t,fragile] \frametitle{Multimodalità}
    \framesubtitle{Speech-to-text: OpenAI con impostazioni di \textit{default}}
    \vspace*{-.7cm}
    {\footnotesize
        \begin{block}{\textit{Controller} di trascrizione}
			{\tiny\inputminted{java}{code/AudioController.java}}
    	\end{block}
        \vspace*{.1cm}
        \begin{itemize}[leftmargin=10pt,align=right]
            \onslide<1->\item[\alert{\faArrowCircleRight}] \alert{\texttt{OpenAiAudioTranscriptionModel}} implementazione di \alert{\texttt{AudioTranscriptionModel}}
            \onslide<2->\item[\alert{\faArrowCircleRight}] Impostazioni di \textit{default}
            \begin{itemize}[leftmargin=10pt,align=right]
                \onslide<3->\item[\alert{\faArrowCircleRight}] Modello Whisper 1
                \item[\alert{\faArrowCircleRight}] \textit{Response format} in JSON
                \item[\alert{\faArrowCircleRight}] Temperatura 0
            \end{itemize}
            \onslide<4->\item[\alert{\faArrowCircleRight}] Formati supportati: \texttt{mp3}, \texttt{mp4}, \texttt{mpeg}, \texttt{mpga}, \texttt{m4a}, \texttt{wav}, \texttt{webm} (documentazione classe \alert{\texttt{AudioTranscriptionPrompt}})
            \onslide<5->\item[\alert{\faExclamationTriangle}] Metodo \alert{\texttt{call}} per allinearsi all'astrazione LLM in Spring AI       \end{itemize}       
    }
\end{frame}    
%
\begin{frame}[t,fragile] \frametitle{Multimodalità}
    \framesubtitle{Speech-to-text: OpenAI con impostazioni \textit{custom} dinamiche}
    \vspace*{-.7cm}   
    {\footnotesize
        \begin{block}{\textit{Controller} di trascrizione}
			{\tiny\inputminted{java}{code/AudioControllerCustom.java}}
    	\end{block}
        \vspace*{.1cm}
        \begin{itemize}[leftmargin=10pt,align=right]
            \onslide<1->\item[\alert{\faArrowCircleRight}] \alert{\texttt{AudioTrascriptionPrompt}} incapsula la risorsa audio e le opzioni di trascrizione (\alert{\texttt{OpenAiAudioTranscriptionOptions}})
            \begin{itemize}[leftmargin=10pt,align=right]
                \onslide<2->\item[\alert{\faArrowCircleRight}] \alert{\texttt{.prompt}} fornisce al LLM contesto riguardo gli argomenti trattati nel \textit{file} audio
                \item[\alert{\faArrowCircleRight}] \alert{\texttt{.language}} evita al LLM l'attività di \textit{language guessing}
                \item[\alert{\faArrowCircleRight}] \alert{\texttt{.responseFormat}} per calibrare il formato di \textit{output} al contesto applicativo di utilizzo
            \end{itemize}
            \onslide<3->\item[\alert{\faExclamationTriangle}] con \texttt{call(audioFile)}, Spring AI si occupa di creare un \textit{bean} \texttt{OpenAiAudioTranscriptionOptions} con i valori di \textit{default}
        \end{itemize}      
    }
\end{frame}    
%
\begin{frame}[t,fragile] \frametitle{Multimodalità}
    \framesubtitle{Speech-to-text: OpenAI con impostazioni \textit{custom} statiche} 
    {\footnotesize
        \begin{block}{Configurazione applicativo}
			{\tiny\inputminted{yaml}{code/application-stt.yml}}
    	\end{block}    
    }
\end{frame}    
%
\begin{frame}[t,fragile] \frametitle{Multimodalità}
    \framesubtitle{Speech-to-text: OpenAI \textit{compatible} API}
    {\footnotesize
        \begin{itemize}[leftmargin=10pt,align=right]
            \onslide<1->\item[\alert{\faExclamationTriangle}] Ollama non supportato
            \onslide<2->\item[\alert{\faExclamationTriangle}] Gemini API supportato \alert{con restrizioni}
            \begin{itemize}[leftmargin=10pt,align=right]
                \onslide<3->\item[\alert{\faArrowCircleRight}] Non possibile utilizzare le astrazioni \alert{\texttt{OpenAiAudioTranscriptionModel}} e \alert{\texttt{AudioTranscriptionPrompt}}
                \onslide<4->\item[\alert{\faArrowCircleRight}] Configurazione statica non ha effetto
                \onslide<5->\item[\alert{\faArrowCircleRight}] Sfruttare modelli con \alert{caratteristiche intrinseche STT}\ldots
                \onslide<6->\item[\alert{\faArrowCircleRight}] \ldots come Gemini 2.5!
            \end{itemize}
        \end{itemize}      
    }
\end{frame}    
%
\begin{frame}[t,fragile] \frametitle{Multimodalità}
    \framesubtitle{Speech-to-text: OpenAI \textit{compatible} API}
    \vspace*{-.7cm}
    {\footnotesize
        \begin{block}{\textit{Upload file} audio su \textit{cloud} temporaneo Google}
			{\tiny\inputminted{bash}{code/gemini-stt.sh}}
    	\end{block}    
    }
\end{frame}    
%
\begin{frame}[t,fragile] \frametitle{Multimodalità}
    \framesubtitle{Speech-to-text: OpenAI \textit{compatible} API}
    {\footnotesize
        \begin{block}{Chiamata Gemini API \textit{generateContent}}
			{\tiny\inputminted{bash}{code/gemini-stt-2.sh}}
    	\end{block}    
        \begin{itemize}[leftmargin=10pt,align=right]
            \onslide<1->\item[\alert{\faArrowCircleRight}] Due strade percorribili
            \begin{itemize}[leftmargin=10pt,align=right]
                \onslide<2->\item[\alertedcircled{1}] Utilizzo Spring Web puro (\texttt{RestTemplate}, \texttt{ResponseEntity}, \texttt{HttpEntity}, \ldots)
                \onslide<3->\item[\alertedcircled{2}] Utilizzo Spring AI (STT attraverso \texttt{ChatClient})
            \end{itemize}
        \end{itemize}         
    }
\end{frame}    
%
\begin{frame}[t,fragile] \frametitle{Progetto Spring AI}
    \framesubtitle{Applicazione e passaggi}
    {\small
    \begin{itemize}[leftmargin=10pt,align=right]
        \onslide<1->\item[\alert{\faArrowCircleRight}] Servizio \textit{speech-to-text} Gemini
        \begin{itemize}[leftmargin=10pt,align=right]
            \onslide<2->\item[\alertedcircled{1}] Creazione interfaccia ed implementazione del servizio
            \onslide<3->\item[\alertedcircled{2}] Modifica controllore MVC
            \onslide<4->\item[\alertedcircled{3}] \textit{Test} delle funzionalità con Postman/Insomnia
        \end{itemize}
    \end{itemize}
    }
\end{frame}
%
\begin{frame}[t,fragile] \frametitle{Progetto Spring AI}
    \framesubtitle{Speech-to-text: Gemini}
        \begin{block}{Interfaccia servizio}
			{\tiny\inputminted{java}{code/MultimodalityService.java}}
    	\end{block}
\end{frame}
%
\begin{frame}[t,fragile] \frametitle{Progetto Spring AI}
    \framesubtitle{Speech-to-text: Gemini}
    \vspace*{-.7cm}
    \begin{block}{Implementazione servizio - I}
		{\tiny\inputminted{java}{code/MultimodalityServiceImpl.java}}
    \end{block}
\end{frame}
%
\begin{frame}[t,fragile] \frametitle{Progetto Spring AI}
    \framesubtitle{Speech-to-text: Gemini}
    \begin{block}{Implementazione servizio - II}
		{\tiny\inputminted{java}{code/MultimodalityServiceImpl-2.java}}
    \end{block}
\end{frame}
%
\begin{frame}[t,fragile] \frametitle{Progetto Spring AI}
    \framesubtitle{Tool calling}
    	\vspace*{-.7cm}
        \begin{block}{Implementazione controllore REST}
			{\tiny\inputminted{java}{code/GeminiMultimodalityController.java}}
    	\end{block}
\end{frame}
%
\begin{frame}[fragile] \frametitle{Codice}
    \framesubtitle{Branch di riferimento}
	\begin{center}
		{\scriptsize \href{https://github.com/simonescannapieco/spring-ai-advanced-dgroove-venis-code.git}{\texttt{https://github.com/simonescannapieco/spring-ai-advanced-dgroove-venis-code.git}}}\\
		\textit{Branch:} \alert{\texttt{16-spring-ai-gemini-ollama-multimodality-stt}}
	\end{center}
\end{frame}